\section{Generative AI：把风险控制前移到全生命周期}

生成式 AI 改变的不只是“产生了新的文件”，还改变了身份伪装、内容变体、应用集成和分发速度。只在最终平台上传处扫描，会把所有责任推到供应链末端。更稳健的做法是把 threat model 分布到 training data、model access、application、hosting/distribution 与 incident response，并在 development、deployment、maintenance 三个阶段持续验证。

\subsection{Threat surface 横跨四层}

本节先建立供应链视角。模型开发者控制训练与 fine-tuning，模型服务控制 access 与 abuse monitoring，应用开发者控制 prompt flow、user identity 和产品 guardrail，分发平台控制 hosting、sharing、reporting 与 removal。不同参与者看到的信号不同，因此不能要求末端平台独自恢复上游已经丢失的 provenance 或账户关系。

\lecturefigure{07-genai-threat-surface.png}{Generative AI 风险横跨开发、模型访问、应用与分发层。}{本地字幕 05:45--10:20、20:30--24:10；Thorn Safety by Design 白皮书；概念重绘。}

\begin{knowledgebox}{读图：为每层写 capability 与 obligation}
Development 层能做 dataset curation、evaluation 与 release gating；model access 层能做 authentication、rate limit、usage monitoring 和 abuse response；application 层能设计 age-appropriate UX、contextual guardrail 与 user reporting；distribution 层能执行 detection、removal、evidence preservation 与 lawful reporting。任何一层都不拥有全部事实，但每层都应保存可向下一层传递的最小安全信号。
\end{knowledgebox}

\begin{warningbox}{不要把 threat model 变成攻击手册}
安全团队需要理解滥用类别、入口和可观测信号，但公开讲义、测试报告与产品文档不应给出可复现的规避步骤、敏感提示词或生成流程。有效的透明度应说明覆盖范围、测试方法、失败类别、响应时间和治理边界，而不是披露能显著降低滥用成本的细节。
\end{warningbox}

\subsection{Safety by Design 是生命周期纪律}

\term{Safety by Design} 指在设计、开发、部署和维护过程中持续纳入风险控制，而不是上线后再加一个 moderation endpoint。它要求团队在定义产品目标时写 threat model，在训练和 release 前做 evaluation/red teaming，在部署时设计 monitoring、escalation 与 user controls，在维护期处理 drift、incident 和 transparency reporting。

\lecturefigure{08-safety-lifecycle.png}{Safety by Design 覆盖 design、development、deployment 与 maintenance，而非一次发布检查。}{本地字幕 20:30--26:20；Thorn/All Tech Is Human principles 与 NIST AI 600-1；概念重绘。}

\begin{knowledgebox}{读图：Release gate 只是中间点}
Design 阶段定义 intended use、misuse 与受影响群体；development 阶段治理数据、模型与测试；deployment 阶段配置 threshold、human review、appeal 和 incident channel；maintenance 阶段监控 drift、更新 policy、修复漏洞并公开进展。一个通过 launch checklist 的模型若没有持续 telemetry 和责任人，几周后仍可能因流量、语言或产品集成变化而失效。
\end{knowledgebox}

\term{red teaming}（红队测试）是在授权、受控环境中主动寻找系统失败模式；它应覆盖模型输出、产品工作流、账户滥用、供应链与 response process。测试结果要转化为 owner、severity、fix、retest 和 residual risk，而不是只展示一组“成功诱导模型”的截图。

\teachervoice{Cordua 把行业承诺描述为一种共同语言：儿童安全专家不可能替每家公司重写整个产品，但可以把反复出现的风险转成可测试的 principles、standards 和 progress reporting，让责任进入正常工程流程。}

\subsection{本章小结}

生成式 AI 需要供应链责任与生命周期控制。末端扫描仍重要，但不能替代上游数据治理、模型访问控制、应用 guardrail、持续测试和透明度。Safety by Design 的价值在于把安全从一次性项目变成可维护系统。

